\section{The source finite grid of short integral harmonics}
\label{sec:sr-source-short-harmonic-grid}

Let \(E\) be a CM field, \(f=[E:\mathbb Q]/2\), and
\(\operatorname{rd}=\operatorname{rd}(E)\).
Give \(E^n\) its standard positive Hermitian form
\(h(x,y)=\sum_kx_k\bar y_k\).
Let \(\Lambda\subset E^n\) be a finitely generated
\(O_E\)-module with
\[
 \Lambda=\{x\in E^n:h(x,\Lambda)\subset O_E\}.
 \tag{HG1}\label{eq:sr-source-harmonic-selfdual}
\]
Take \(p+q=r\le n\), finite \(T\), and minor exponents
\(c,c'\) as in Section~\ref{sec:sr-source-harmonic-orbit};
write \(D=\mathrm{typeDegree}(c,c')\).
The construction below uses one finite set of lattice-vector
arrays for all coefficient fields and all embeddings in its
statement.

\begin{srproposition}\label{prop:sr-source-short-harmonic-grid}
If \(D\ge1\), put \(d_g=(2r+1)D\) and
\[
 \rho_0=2f^{3/2}n^2\operatorname{rd},\qquad
 R=(2fn)d_g\rho_0,\qquad
 R_*=(r+1)!\,R^{2r+1}.
\]
There is a finite set \(S\) of pairs
\((u_1,\ldots,u_p;w_1,\ldots,w_q)\) in
\(\Lambda^p\times\Lambda^q\), with \(h(u_a,w_b)=0\),
such that
\[
 |S|\le(d_g+1)^{2fn(2p+q)}.
 \tag{HG2}\label{eq:sr-source-harmonic-grid-card}
\]
For any characteristic-zero field \(K\) and maps
\(\tau_t:E\to K\) for which the \(2|T|\) maps
\(\tau_t,\tau_t\circ\mathrm{conj}\) are pairwise distinct,
\[
 \operatorname{span}_K
 \left\{{\cal M}_{c,c'}\bigl((\tau_tu_a)_{t,a},
                         (\tau_t\bar w_b)_{t,b}\bigr):
                         (u,w)\in S\right\}
            ={\cal H}_K(c,c').
 \tag{HG3}\label{eq:sr-source-harmonic-grid-span}
\]
Each displayed polynomial is harmonic and lies in the
subring generated by the \(\mathbb Z\)-integral constants
of \(K\), the forms \(X_{t,a}(\tau_t v)\), and the forms
\(Y_{t,b}(\tau_t\bar v)\), with \(v\in\Lambda\).
This is the source's lattice-generated coefficient ring.
For every choice of complex embeddings \(\tau_t\), without
a distinctness requirement for this estimate, its
coefficient \(\ell^1\)-norm is at most
\[
       (r\sqrt n\,R_*)^D.
 \tag{HG4}\label{eq:sr-source-harmonic-grid-l1}
\]
\end{srproposition}

\begin{proof}
The actual short-vector producer gives \(m=2fn\) vectors
\(v_i\in\bar\Lambda\) spanning \(E^n\) over \(\mathbb Q\),
indeed a \(\mathbb Q\)-basis, such that
\[
 \left(\sum_k|\sigma(v_{i,k})|^2\right)^{1/2}
          \le\rho_0\quad\hbox{for every }\sigma:E\to\mathbb C.
 \tag{HG5}\label{eq:sr-source-short-vector-basis-use}
\]
This is the \(2fn\)-vector result proved from the trace-dual
minimum and Euclidean transference in
Section~\ref{sec:sr-source-gaussian-sup}.
It applies to \(\bar\Lambda\) because conjugation is a
semilinear involution, preserves finite generation, and
\(h(\bar x,\bar y)=h(y,x)\); consequently
\(\bar\Lambda\) also satisfies (HG1).

Choose arrays of integers
\(0\le A_{a,i},A'_{a,i},B_{b,i}\le d_g\), and set
\[
 e_a=\sum_iA_{a,i}v_i,\qquad
 e'_a=\sum_iA'_{a,i}v_i,\qquad
 g_b=\sum_iB_{b,i}v_i .
\]
Let \(M_{a',a}=h(e'_{a'},e_a)\), and form the source
adjugate orthogonalizer
\[
 f_b=\det(M)g_b-
       \sum_{a,a'}\operatorname{adj}(M)_{a,a'}
                              h(g_b,e_a)e'_{a'}.
 \tag{HG6}\label{eq:sr-source-adjugate-orthogonalizer}
\]
All Gram entries are in \(O_E\), so \(f_b\in\bar\Lambda\),
with no division by \(\det M\).
The identity \(\operatorname{adj}(M)M=\det(M)I\) gives
\(h(f_b,e_a)=0\), even when \(M\) is singular.
Define \(u_a=\bar e_a\), \(w_b=\bar f_b\), and let \(S\)
be the image of this integer grid.
Then \(u_a,w_b\in\Lambda\) and \(h(u_a,w_b)=0\).
There are \(m(2p+q)\) grid entries, proving (HG2).
The polynomials in (HG3) are exactly
\({\cal M}_{c,c'}((\tau_t\bar e_a),(\tau_tf_b))\).
This conjugation step is necessary for the actual theta
letters, whose second block uses the conjugate lattice
row form.

Here is the spanning proof, including the simultaneous
embedding assertion.  Suppose \(v_i\) span \(E^n\) over
\(\mathbb Q\), and let \(\sigma_j:E\to K\) be distinct
field maps.  The \(K\)-linear map
\[
 K^m\longrightarrow K^{J\times n},\qquad
 (z_i)\longmapsto
       \left(\sum_i z_i\sigma_j(v_{i,k})\right)_{j,k}
 \tag{HG7}\label{eq:sr-source-embedding-surjection}
\]
is onto.
If a linear functional annihilates its range, it also
annihilates the image of every vector of \(E^n\), by
\(\mathbb Q\)-spanning.
Apply it to a vector with just one nonzero coordinate
\(x\in E\).  This gives
\(\sum_j a_{j,k}\sigma_j(x)=0\) for all \(x\).
Distinct field maps are linearly independent:
a shortest nontrivial relation, evaluated at \(xy\)
and subtracted from \(\sigma_{j_0}(y)\) times the
relation at \(x\), has fewer terms after choosing
\(y\) separating two maps, a contradiction.
Hence every \(a_{j,k}=0\).  The dual annihilator is
zero, which proves (HG7).

Now replace the integer arrays \(A,A',B\) temporarily
by independent \(K\)-parameters, and use the bilinear dot pairings whose coefficients
are obtained after applying \(\tau_t\).
The first coefficient array is
\(\epsilon_{t,a}=\sum_i A_{a,i}\tau_t(\bar v_i)\);
the other two arrays use \(\tau_t(v_i)\).
Surjectivity (HG7), applied to
\(\sigma_j=\tau_t,\tau_t\circ\mathrm{conj}\), lets
one prescribe both embedding arrays independently
for each row.
Given an independent perpendicular pair
\((\epsilon_{t,a},\phi_{t,b})\), prescribe the first
array to be \(\epsilon\), choose an auxiliary array
\(\epsilon'\) with
\(\epsilon'_{t,a'}\cdot\epsilon_{t,a}=\delta_{a',a}\),
and prescribe the third array to be \(\phi\).
For these choices \(M=I\), and (HG6) is precisely
\(\phi\).  Thus the parameterized minor products
contain every orbit member of \(P_0\).
Conversely every parameter value is a perpendicular
pair by the adjugate identity, so its minor product
belongs to the orbit span, including dependent pairs
by the flag argument at the end of
Section~\ref{sec:sr-source-harmonic-orbit}.
The span of all parameter values is therefore exactly
\({\cal H}_K(c,c')\).

Finally this same span is obtained on the finite grid.
Every entry of \(M\) is homogeneous of parameter
degree \(2\), every adjugate entry has degree
\(2(p-1)\) when present, and every entry of \(f_b\)
has degree \(2p+1\).
Thus, after exchanging the coefficient and letter
variables, the vector-valued polynomial giving the
minor product has total parameter degree
\[
 \sum_t\left(\sum_i i c_{t,i}
                 +(2p+1)\sum_j j c'_{t,j}\right)
       \le(2p+1)D\le d_g.
\]
For \(p=0\), (HG6) is simply \(f_b=g_b\), with
\(\det(\varnothing)=1\), so the same degree assertion
holds.
A scalar polynomial of degree at most \(d_g\) in
each variable vanishing on
\(\{0,\ldots,d_g\}^{m(2p+q)}\) is zero: apply the
one-variable root bound successively.
Apply that assertion to every linear functional
annihilating the span of the grid values.
It annihilates every parameter value, proving that
the grid span is the full parameter span.
This proves (HG3), with one grid before \(K,\tau\)
are chosen.

We prove the estimate last, with the source's constants.
In each complex embedding,
\(\|e_a\|,\|e'_a\|,\|g_b\|\le m d_g\rho_0=R\).
Here \(r,n,f,D\ge1\), \(\operatorname{rd}\ge1\), so
\(R\ge1\).
The Gram entries have absolute value at most \(R^2\).
Leibniz expansion in (HG6) gives
\[
 \|f_b\|
 \le \{p!+p^2(p-1)!\}R^{2p+1}
       =(p+1)!\,R^{2p+1}\le R_*
\]
for \(p\ge1\); for \(p=0\), \(\|f_b\|\le R\le R_*\).
Every linear row form in the minor product has
coefficient \(\ell^1\)-norm at most \(\sqrt n R_*\)
by Cauchy--Schwarz.
A leading determinant of size \(i\le r\) consequently
has \(\ell^1\)-norm at most
\(i!(\sqrt n R_*)^i\le(r\sqrt n R_*)^i\).
Submultiplicativity and (HO1) prove (HG4).
Products of lattice row forms have the asserted
integral-algebra membership, and their contractions
vanish because their two coefficient arrays are
perpendicular.
\end{proof}

\begin{srlemma}\label{lem:sr-source-short-harmonic-log}
Suppose \(L\ge2\),
\[
 \log\operatorname{rd}\le2L,\quad \log f\le2L,\quad
 \log(1+D)\le2L,\quad \log n\le L .
\]
The polynomials of Proposition~\ref{prop:sr-source-short-harmonic-grid}
satisfy
\[
       \log\|Q\|_1\le140rLD.
 \tag{HG8}\label{eq:sr-source-short-harmonic-log}
\]
The assertion is valid also at \(\|Q\|_1=0\), with the
source convention \(\log0=0\).
\end{srlemma}
\begin{proof}
Put \(u=e^L\ge3\).
The hypotheses and \(1\le r\le n\) give
\[
 n,r,\sqrt n\le u,\quad f,\operatorname{rd},D,
                       2r+1,r+1\le u^2,\quad
 f^{3/2}\le u^4,\quad (r+1)!\le u^{4r}.
\]
In particular
\((2fn)d_g\rho_0\le u^{17}\), and
\[
 r\sqrt n\,R_*
       \le u^{\,2+4r+17(2r+1)}\le u^{140r}.
\]
Use (HG4) and take logarithms when \(Q\ne0\).
For \(Q=0\), the right side is nonnegative.
\end{proof}

\subsection{The literal theta letters and all coefficient embeddings}
Let \(H\) be the actual host, \(D_w\) its datum,
\(p=D_w.p\), \(q=D_w.q\), \(F=c.F\) the actual coefficient
field, and \(a(\tau,i)\) the source exponent array.
The source's letter bijection is
\[
 (\tau,i,k)\longmapsto
 \begin{cases}
   (\tau,x,i,k),&\tau\in\Psi,\\
   (\bar\tau,y,i,k),&\tau\notin\Psi .
 \end{cases}
 \tag{HG9}\label{eq:sr-source-theta-letter-bijection}
\]
The row indices here retain their order within their
block; the reversal of negative columns in \(P_0\)
is a different operation.
For the lattice letter
\(\ell_{\tau,i}(v)=\sum_k\mathrm{embF}_c(\tau)(v_k)
\,Z_{\tau,i,k}\), the bijection gives
\(X_{\tau,i}(\tau v)\) on \(\Psi\) and
\(Y_{\bar\tau,i}(\bar\tau(\bar v))\) off \(\Psi\).
It therefore transports the actual lattice-generated
ring \(\mathrm{symO}(c,\Lambda)\), as well as every
Laplacian and row-addition equation, to the rings and
equations just used.

If the actual harmonic space \(\mathrm{HHa}(a)\)
is nonzero, every exponent row is antitone.
Indeed its row-raising derivations annihilate a
nonzero element, and its row-diagonal derivations
have weights \(a(\tau,i)\); the same finite-degree
\(\mathfrak{sl}_2\) calculation used in
Section~\ref{sec:sr-source-harmonic-orbit} proves
\(a(\tau,i)\ge a(\tau,j)\) for \(i<j\).
Put
\[
 c_{\tau,i}=a(\tau,i)-a(\tau,i+1),\qquad
 c'_{\tau,j}=a(\bar\tau,j)-a(\bar\tau,j+1),
\]
using zero after the last row.
Then the letter bijection identifies
\(\mathrm{HHa}(a)\) with
\({\cal H}_F(c,c')\), and
\(\mathrm{typeDegree}(c,c')=\mathrm{DF}(D_w,a)\).
If the exponent condition fails, the harmonic space
is zero, and we take the empty family.

The maps \(\mathrm{embF}_c(\tau)\) and their conjugates,
for \(\tau\in\Psi\), are distinct.
After inclusion in \(\mathbb C\), equality would
identify two field embeddings, and a CM-type
embedding cannot equal the conjugate of another
one in the same CM type.  This last fact follows
already from \(\tau(H.\mathrm{rt})=H.\mathrm{rt0}\):
the conjugate has the opposite nonzero imaginary
part.  Thus (HG3) applies over the actual \(F\).

For any embedding \(\sigma:F\to\mathbb C\),
the images of the coefficients are obtained by the
maps \(\sigma\circ\mathrm{embF}_c(\tau)\).
The norm estimate (HG4) applies to every such tuple
of embeddings.  Renaming the variables is bijective
and preserves the coefficient \(\ell^1\)-norm.
Consequently the finite family, transported back
by (HG9), consists of elements of
\(\mathrm{harmL}(c,\Lambda,a)\), spans the whole
\(\mathrm{HHa}(a)\) over \(F\), satisfies
\[
 \log\sum_m|\sigma(Q_m)|\le
          140rL\,\mathrm{DF}(D_w,a)\quad(\sigma:F\to\mathbb C),
 \tag{HG10}\label{eq:sr-source-harmL-all-embeddings}
\]
and has cardinality at most
\[
 \bigl((2r+1)\mathrm{DF}(D_w,a)+1\bigr)^{2fn(2p+q)}.
\]
If \(\mathrm{DF}(D_w,a)=0\), every member of
\(\mathrm{HHa}(a)\) is constant.
Use \(\{1\}\) if \(1\) belongs to that space, and the
empty family otherwise; every displayed bound remains
valid.  This handles the source's zero-degree cases
separately from Proposition~\ref{prop:sr-source-short-harmonic-grid}.

\subsection{The exact \(20000\) harmonic bound on the host window}
Use the full \(\mathrm{CovGmin.Hyp}(D,T,j,Y,e,o,L_b,a,n)\)
tuple of Section~\ref{sec:sr-source-raw-sup-9000},
including its actual good coefficient field.
Write \(F=Y.d.c_2.F\), \(f=D.R.f\), \(r=D.R.r\),
\(d=\mathrm{DF}(T.Q.Dw(j),a)\), and
\(\kappa=\mathrm{sizeK}(T.Q.Dw(j),a,n)\).
Let \(\Lambda\) be unimodular and trivial at the actual
integer \(\mathrm{badC}(Y)\).
It is finitely generated, as required above:
the triviality witness supplies an integer \(k\)
coprime to \(\mathrm{badC}(Y)\) with
\(k\Lambda\subset O_E^n\).
Since \(2\mid\mathrm{badC}(Y)\), \(k\ne0\);
injective multiplication by \(k\), and Noetherianity
of \(O_E^n\), give finite generation.
Unimodularity gives (HG1).
The host rank identity gives \(p+q=r\le n\).

\begin{srproposition}\label{prop:sr-source-harmonic-20000}
Every \(P\in\mathrm{harmL}(Y.d.c_2,\Lambda,a)\) belongs
to the \(F\)-span of those \(Q\) in that same space with
\[
 \mathrm{polyNorm}_{Y.d.c_2}(Q)
       \le \exp\!\left(20000(frL_b)^2\kappa\right).
 \tag{HG11}\label{eq:sr-source-harmonic-20000}
\]
The constant \(20000\) is chosen before the host,
window, coefficient field, weight and lattice.
\end{srproposition}
\begin{proof}
The strict positive-row condition, evenness of \(n\),
and \(p\ge1\) imply \(\kappa\ge1\) and \(n\le2\kappa\):
choose a CM-type embedding and its first positive row;
that row is an integer strictly greater than \(n/2\)
and its absolute value is one summand of \(\kappa\).
Each exponent is at most the absolute value of its
corresponding weight row, so, summing over all rows,
\[
                d\le [E:\mathbb Q]r\kappa=2fr\kappa.
\]
Take the source value
\[
 L=2+3L_b+\log(1+d)+\log n.
\]
It is at least \(2\), and the field-size bound gives
\(\log\operatorname{rd},\log f\le3L_b\le2L\).
The other two hypotheses of
Lemma~\ref{lem:sr-source-short-harmonic-log} follow
from the definition of \(L\).
Moreover \(1+d\le4fr\kappa\), since \(fr\kappa\ge1\).
Using \(\log(fr)\le L_b\) and
\(\log\kappa\le20L_b\) gives
\[
 \log(1+d)\le3+21L_b,\qquad
 \log n\le1+20L_b,\qquad L\le50L_b .
\]
For the finite family in (HG10), set
\(B=\exp(140rLd)\).
The literal coefficient norm in the source is
\[
 \mathrm{polyNorm}_c(Q)^2
       =\sum_{\sigma:F\to\mathbb C}\sum_m|\sigma(Q_m)|^2.
\]
Since each \(\ell^2\)-norm is at most its
\(\ell^1\)-norm,
\(\mathrm{polyNorm}_c(Q)\le\sqrt{[F:\mathbb Q]}B\).
The retained \(\mathrm{Good}\) witness gives
\(\log[F:\mathbb Q]\le o\), hence
\[
 \log\mathrm{polyNorm}_c(Q)
             \le o/2+140rLd
             \le20000(frL_b)^2\kappa
\]
when this norm is positive; the zero case is immediate.
For the last inequality, use
\[
 \begin{split}
 140rLd&\le14000(fr^2L_b)\kappa
                      \le14000(frL_b)^2\kappa,\\
 o/2&\le(391/2)(frL_b)^2
                      \le(391/2)(frL_b)^2\kappa .
 \end{split}
\]
Here \(f,L_b,\kappa\ge1\), so each comparison has
the stated direction.
Thus the whole finite spanning family satisfies
(HG11), which proves the proposition.
\end{proof}

\paragraph{Frozen-source correspondence.}
The following files have prefix
\srcid{src/Development/AbelianVarietyTheory/Proofs/}.
Under \srcid{CuspFormEvaluations/}:
\srcid{SpanGridPerpVecEqOrbitSpanOfLe.lean}
(\srcid{43867edccdfb7b1c14649d22af2cd20e87731621});
\srcid{ExistsShortHarmonicFamilyLinear.lean}
(\srcid{175e84697ec07b92f83460d2115d685af0a91693});
\srcid{ExistsShortHarmonicFamilyFewColumnsLinear.lean}
(\srcid{e40b0f4084981ebccad52f0692b3dd45ca015888});
\srcid{L1NormRowBlockLeadingMinorProductLe.lean}
(\srcid{8498ec8435a193542b52c07806b1951ab9a1fd9a});
\srcid{LogShortFamilyNumbersLe.lean}
(\srcid{2f96394d1cb1e84f2c37b09ceaccd153ea02d249});
\srcid{ThetaLetterEquivConjugate.lean}
(\srcid{586c5a03c6e618eca3d06c7f4ebe8793994b7d19});
\srcid{ExistsShortHarmonicFamilyUnimodular.lean}
(\srcid{a3e3ffcb0a590b3e29bc36ba72dd2dd4222e7841}).
The finite grid interpolation is
\srcid{CovolumeBounds/EqZeroOfEvalGridEqZeroOfTotalDegreeLe.lean}
(\srcid{ac7acd584a4e382a387d99bd0653a8bae70f32b2}).
The final window estimate is the declaration
\srcid{CovGmin.PrismSpan.harm_short} in
\srcid{DivisorProofs/PolynomialCoefficientEmbeddingNormLeOfL1Bound.lean}
(\srcid{c13735d04f35ab319d44643cd3bc173a04117d7b}).
The result is a proved harmonic supplier.  The actual passage
from these polynomials to the same all-test-ring cusp sections
still requires the separately identified uniformization,
short-rule and stable theta-spanning producers; those have
not been replaced by an abstract existence assertion here.
